\chapter{The damped frame also enlarges the pure-testing state chart}
\label{component:p8cfada778490}
\noindent\textbf{Component }\path{p8cfada778490}\par
{\small\noindent Source: \path{R17/proofs/c1_small_eta/damped_state_extension.tex}\par}
{\small\noindent SHA256: \path{8cfada778490ddc7e1f29ebd9c7ef1007bb5c294c670c2f2a38d491b1309e5dd}\par}
{\small\noindent Inventory use: DS2--DS3 | C14\_\allowbreak{}ENTIRE\_\allowbreak{}ORIGINAL\_\allowbreak{}PARENT\_\allowbreak{}GRAPH; precise subordinate uses retained below\par}
\medskip
\section{The damped frame also enlarges the pure-testing state chart}

\begin{proposition}
On the entire unchanged C14.3b source, with no condition on $x/\eta$,
the complete local parent is tensor SOS on
\[
 Z\ne0,\qquad X/Z\text{ bounded},\qquad
                         u^2W/Z\text{ bounded}.       \tag{DS1}
\]
No bound on $Y/Z$ is imposed. The literal $Z=0$ endpoint is already
covered by the direct sign-core theorem. This extends R16's local
bounded-$W/Z,X/Z$ state chart, since $u$ is infinitesimal.
\end{proposition}

Use the exact damped identity DG3 with $s=x+\eta$. Its strictly
positive source frame and all residual allocations except DG10 are
valid without the source restriction $x/\eta$ bounded. Thus only
$\epsilon uvQ\xi$ requires a new allocation. Put
$\phi=\xi/Z=X/Z+hu$, which is bounded on DS1. There is the exact
identity
\[
 \boxed{\quad
 \epsilon uvQ\xi
 =\epsilon U(W/Z)E\xi
       -\epsilon u\phi E^2/\eta-\epsilon uvP\xi.
 \quad}                                                   \tag{DS2}
\]
It follows from $Q=E/\eta-P$ and $vZ=\eta uW-E$; its only testing
inverse is $Z^{-1}$.

On the DG8 source frame, $P$ has a bounded allocation and
$\sqrt{\eta s}\,Q$ has a bounded allocation with unit coefficient.
Use $E=\eta P+\eta Q$ and $\epsilon=O(s)$. The first term of DS2
has coefficients $O(\eta s)U(W/Z)$ and
$O(\sqrt{\eta s})U(W/Z)$ on $PN_x$ and $N_QN_x$, respectively;
these are infinitesimal under DS1. In the second term, the three
coefficients after expanding $E^2/\eta$ are
\[
 O(\eta s)u\phi,\qquad O(\sqrt{\eta s})u\phi,
                                      O(1)u\phi
\]
on $P^2$, $PN_Q$, and $N_Q^2$. They are all infinitesimal, since
$u$ is infinitesimal and $\phi$ bounded. The last term of DS2 has
coefficient $O(suv)$ on $PN_x$, also infinitesimal. This allocates
the complete missing normal on the existing strict frame; no old
$E/\sqrt\eta$ pivot is reintroduced. The same retained positive
rows and sign-core join used in DG3 prove the assertion.

\paragraph{A direct state exit and the resulting frontier.}
For $Z\ne0$, if $X/Z$ is unbounded, then
$(X+uZ)/Z$ is unbounded. The following portion of the original signed
half core is strictly positive in the pure testing field:
\[
 \frac{1+v^2}{2}\big\{(X+uZ)^2-(1+u^2)Z^2\big\}>0.
                                                               \tag{DS3}
\]
Every remaining original packet has a nonnegative pure-source
coefficient, including the other two half-core squares. A real
closed testing field therefore gives a direct tensor-SOS certificate.
This exit uses no contact or source estimate. Consequently the
remaining nontrivial real state chart, after the sign gauges, can be
restricted to
\[
 WZ>0,\quad Z\ne0,\quad X/Z\text{ bounded},\quad
                              u^2W/Z\text{ unbounded}.         \tag{DS4}
\]
After normalizing a $W$-unit state by $W=1$, DS4 says
$Z/u^2$ is positive infinitesimal. In addition, the separate DG
source admission removes every source with $x/\eta$ bounded. These
are pure-source and pure-testing reductions, not a comparison of
their rates.

For a complex state satisfying the modulus version of DS1, use a
common testing phase to make nonzero $Z$ real. Its real part satisfies
DS1 and its imaginary part has $Z=0$; each has the indicated complete
certificate. All displayed formulas include the local axes by their
literal values. The focused checker verifies DS2 and the exact
original-parent packet accounting used in DS3. No all-tests claim is
made for the remaining source chart $x/\eta$ unbounded.

\paragraph{Precise remaining normalized corner and residual.}
Combine DG1, DS3--DS4, the inherited reciprocal-angle local theorem,
and the negative-core consequence $uvW/Z$ bounded from this
revision's sign-core Schur formula. On a normalized $W$-unit state,
the only remaining C14.3b corner has
\[
 \begin{gathered}
 x/\eta\text{ unbounded},\qquad
 u,v>0\text{ infinitesimal},\quad v/u\text{ infinitesimal},\\
 W=1,\quad Z>0,\quad Z/u^2\text{ infinitesimal},\quad
 X/Z\text{ bounded},\quad uv/Z\text{ bounded}.
 \end{gathered}                                             \tag{DS5}
\]
All original source openings and $\beta\ge0$ remain. The two last
terms of DS2 are already infinitesimal Grams on the strict damped
frame throughout DS5. The sole unallocated expression is exactly
\[
                         \epsilon U(W/Z)E\xi.               \tag{DS6}
\]
As $s\asymp x$ on this pure-source branch, its coefficients after
$E=\eta P+\eta Q$ have orders
$\eta x\,UW/Z$ and $\sqrt{\eta x}\,UW/Z$ on $PN_x$ and
$N_QN_x$. The negative-core bound on $uvW/Z$ only rewrites the second
as $O(\sqrt{\eta x}\,u/v)$; it does not bound it in independent
source/testing fields. Thus DS5 is a genuine remaining chart and
DS6 a precise remaining complete normal. No source/testing ordering
or additional strict pivot is claimed to remove it.
